\documentclass[11pt]{article}
\usepackage{amsmath,amssymb,amsthm}
\usepackage[margin=1in]{geometry}

\newtheorem{theorem}{Theorem}
\newtheorem{lemma}[theorem]{Lemma}
\newtheorem{corollary}[theorem]{Corollary}
\newtheorem{proposition}[theorem]{Proposition}
\theoremstyle{definition}
\newtheorem{definition}[theorem]{Definition}
\theoremstyle{remark}
\newtheorem{remark}[theorem]{Remark}

\newcommand{\orb}{\operatorname{orb}}
\newcommand{\supp}{\operatorname{supp}}
\newcommand{\des}{\operatorname{des}}
\newcommand{\maj}{\operatorname{maj}}
\newcommand{\Aalt}{A^{\mathrm{alt}}}
\newcommand{\thetaalt}{\theta^{\mathrm{alt}}}

\title{Sharp Obstruction to the Bigraded Coset Poincar\'e Identity via Adjoint Pairing}
\author{Clio}
\date{2026-07-29}

\begin{document}
\maketitle

\begin{abstract}
Let $\mu \vdash n$, $P_\mu(x;t) = \sum_\gamma c_\gamma(t) \Aalt_\gamma$ be
the expansion in alternating Demazure atoms. Let
$A_\mu(t,q) = \sum_{w \in W_\mu} t^{\des(w)} q^{\maj(w)}$ be the bi-Mahonian
generating function over rearrangements of the multiset $\{\mu_1,\ldots,\mu_n\}$.
The proof of $c_\mu(t) = A_\mu(1,t)$ [Clio 2026-07-28] uses the linear
functional $L(f) = \sum_\sigma t^{\ell(\sigma)}[x^\sigma]f$ and an
adjoint-pairing argument. It is natural to conjecture a bigraded lift
\[
\text{(★)} \qquad L^{(t,q)}(\Aalt_\gamma) \;\stackrel{?}{=}\;
   \delta_{\gamma,\mu}\,t^{\des(\gamma)} q^{\maj(\gamma)},
\qquad L^{(t,q)}(f) := \sum_\sigma t^{\des(\sigma)} q^{\maj(\sigma)}[x^\sigma]f,
\]
which would give $c_\mu \cdot t^{\des(\mu)} q^{\maj(\mu)} = A_\mu(t,q)$, a
genuine two-parameter refinement.

We prove two things:

\begin{enumerate}
\item (★) is \emph{false}: for every $\mu$ with $|\orb(\mu)| \ge 2$, there exist
$\gamma \ne \mu$ with $L^{(t,q)}(\Aalt_\gamma) \ne 0$.

\item (\emph{Rigidity Theorem}) Any weight function $w: \orb(\mu) \to R$
satisfying the descent-swap condition (Definition~\ref{def:ds}) is forced
to be $w(\sigma) = C \cdot t^{\ell(\sigma)}$ for a constant $C$.
\end{enumerate}

The rigidity theorem shows that no bigraded weight
$w: \orb(\mu) \to \mathbb Z[t,q]$ can drive the adjoint-pairing method of
[Clio 2026-07-28] to a genuinely $q$-graded refinement: the machinery is
saturated by the choice $w(\sigma) = t^{\ell(\sigma)}$. A bigraded lift, if
one exists, must come from a strictly different route.
\end{abstract}

\section{Setup}

Recall from [Clio 2026-07-28] the setting: $\mu = (\mu_1 \ge \cdots \ge \mu_n \ge 0) \vdash n$,
$\orb(\mu) = \{\text{distinct rearrangements of } \mu\}$, $\ell(\gamma)$
the Bruhat length of $\gamma$ relative to $\mu$ (number of adjacent
transpositions in a shortest bubble word from $\mu$ to $\gamma$), and
$\Aalt_\gamma = \thetaalt_{i_k}\cdots \thetaalt_{i_1}(x^\mu)$ with
$\thetaalt_i(f) = \frac{x_{i+1} - t x_i}{x_i - x_{i+1}}(f - s_i f)$.
For $\sigma \in \orb(\mu)$, viewed as a word of length $n$, define
$\des(\sigma) = |\{i : \sigma_i > \sigma_{i+1}\}|$ and
$\maj(\sigma) = \sum_{i : \sigma_i > \sigma_{i+1}} i$.

\section{Numerical falsification of (★)}\label{sec:falsify}

\begin{proposition}\label{prop:false}
For $\mu = (2,2,0)$ and $\gamma = (2,0,2) \in \orb(\mu)$,
\[
   L^{(t,q)}(\Aalt_{\gamma}) \;=\; qt\,(1 - qt) \;\ne\; 0.
\]
In particular, (★) is false: $\gamma \ne \mu$ but $L^{(t,q)}(\Aalt_\gamma) \ne 0$.
\end{proposition}

\begin{proof}
Direct computation. From [Clio 2026-07-28, Sec.\ 3.2],
\[
  \Aalt_{(2,0,2)} \;=\; -t\, x_1^2 x_2^2 \;+\; (1-t)\, x_1^2 x_2 x_3 \;+\; x_1^2 x_3^2.
\]
The orbit is $\orb(\mu) = \{(2,2,0),(2,0,2),(0,2,2)\}$ with $(\des,\maj)$-statistics
$(1,2), (1,1), (0,0)$ respectively. Reading off coefficients from
$\Aalt_{(2,0,2)}$: $[x^{(2,2,0)}] = -t$, $[x^{(2,0,2)}] = 1$, $[x^{(0,2,2)}] = 0$.
Hence
\[
  L^{(t,q)}(\Aalt_{(2,0,2)})
    \;=\; t^{1}q^{2}(-t) \;+\; t^{1}q^{1}(1) \;+\; t^{0}q^{0}(0)
    \;=\; qt - q^2 t^2 \;=\; qt(1-qt).
\]
Since this is nonzero as a polynomial in $t,q$, (★) fails.
\end{proof}

\begin{remark}\label{rem:general}
The failure is generic. A SymPy probe confirms that for
$\mu \in \{(2,1,0), (2,2,0), (2,2,1), (3,2,1), (2,2,0,0)\}$, \emph{every}
off-diagonal $\gamma \ne \mu$ has $L^{(t,q)}(\Aalt_\gamma) \ne 0$. Full
per-$\gamma$ output at
\texttt{\char`\~/projects/probes/2026-07-29-bigraded-coset-poincare/probe1.out}.
\end{remark}

\begin{remark}[Sanity: the identity $\sum c_\gamma(t) L^{(t,q)}(\Aalt_\gamma) = A_\mu(t,q)$ holds]
Applying $L^{(t,q)}$ to $P_\mu = \sum c_\gamma \Aalt_\gamma$ gives
$\sum c_\gamma(t) L^{(t,q)}(\Aalt_\gamma) = L^{(t,q)}(P_\mu) = A_\mu(t,q)$
(the last equality because $[x^\sigma] P_\mu = 1$ for $\sigma \in \orb(\mu)$
by Macdonald III.\S 2). So $A_\mu(t,q)$ \emph{does} equal a weighted sum of
$L^{(t,q)}(\Aalt_\gamma)$ over $\gamma$, but the individual terms conspire
non-diagonally.
\end{remark}

\section{The Rigidity Theorem}

We now identify the sharp structural obstruction.

\begin{definition}\label{def:ds}
Let $R$ be a commutative ring containing an indeterminate $t$. A function
$w: \orb(\mu) \to R$ satisfies the \emph{descent-swap condition} if
\[
  w(s_i \sigma) \;=\; t \cdot w(\sigma)
  \qquad\text{whenever } \sigma_i > \sigma_{i+1}.
\]
\end{definition}

\begin{theorem}[Rigidity of the adjoint-pairing weight]\label{thm:rigidity}
Every function $w: \orb(\mu) \to R$ satisfying the descent-swap condition
has the form
\[
  w(\sigma) \;=\; w(\mu) \cdot t^{\ell(\sigma)}
  \qquad \text{for all } \sigma \in \orb(\mu).
\]
In particular, up to a global scalar in $R$, the only such weight function
is $\sigma \mapsto t^{\ell(\sigma)}$.
\end{theorem}

\begin{proof}
Let $\sigma \in \orb(\mu)$ and let $(i_1, \ldots, i_k)$ with $k = \ell(\sigma)$
be a Bruhat-minimum bubble word taking $\mu$ to $\sigma$; that is, setting
$\eta_0 = \mu$ and $\eta_j = s_{i_j} \eta_{j-1}$, we have $\eta_k = \sigma$
and $\ell(\eta_j) = j$ for each $j = 0, 1, \ldots, k$.

At each step $j$: since $\ell(\eta_j) = \ell(\eta_{j-1}) + 1$, the transposition
$s_{i_j}$ acts on $\eta_{j-1}$ by \emph{creating} an inversion. Equivalently
(swapping positions $i_j$ and $i_j+1$ of $\eta_{j-1}$ raises the Bruhat length
iff the entry at position $i_j$ is $>$ the entry at position $i_j+1$),
$\eta_{j-1}$ has a descent at $(i_j, i_j+1)$: $(\eta_{j-1})_{i_j} > (\eta_{j-1})_{i_j+1}$.
Hence the descent-swap condition applies to $\eta_{j-1}$ at position $i_j$:
\[
  w(\eta_j) \;=\; w(s_{i_j} \eta_{j-1}) \;=\; t \cdot w(\eta_{j-1}).
\]
Iterating,
\[
  w(\sigma) \;=\; w(\eta_k) \;=\; t \cdot w(\eta_{k-1}) \;=\; \cdots
                 \;=\; t^k \cdot w(\eta_0) \;=\; t^{\ell(\sigma)} \cdot w(\mu).
\]
This holds for every $\sigma$, completing the proof.
\end{proof}

\begin{corollary}[No genuine $q$-refinement via adjoint pairing]\label{cor:no-q}
Let $w: \orb(\mu) \to \mathbb Z[t,q]$ be any bigraded weight function, and
consider the linear functional
\[
   L_w(f) \;:=\; \sum_{\sigma \in \orb(\mu)} w(\sigma) [x^\sigma] f.
\]
Assume that Step (a) of [Clio 2026-07-28, proof of Lemma~5] (cancellation of
the orbit-supported part of $\widetilde{\thetaalt_i}(Q_w)$ where
$Q_w := \sum_\sigma w(\sigma) x^\sigma$) is required to hold, in the sense
that the coefficient of $x^\tau$ in $\widetilde{\thetaalt_i}(Q_w)$
restricted to $\tau \in \orb(\mu)$ vanishes for every $i, \tau$. Then
\[
   w(\sigma) \;=\; w(\mu) \cdot t^{\ell(\sigma)}.
\]
In particular, $w$ has no dependence on $\sigma$ beyond $t^{\ell(\sigma)}$;
any $q$-dependence is a global multiplicative factor $w(\mu) \in \mathbb Z[t,q]$
and hence contributes only a global rescaling of $L_w$.
\end{corollary}

\begin{proof}
Reading off the diagonal-and-swap contributions from
[Clio 2026-07-28, eqs.\ (2)--(3)] and grouping the descent--ascent pair
$(\sigma, s_i\sigma)$ with $\sigma$ having a descent at $(i,i+1)$: the
coefficient of $x^\sigma$ in $\widetilde{\thetaalt_i}(Q_w)$ receives
$w(\sigma)(-t) + w(s_i\sigma)(1)$, and the coefficient of $x^{s_i\sigma}$
receives $w(\sigma)(t) + w(s_i\sigma)(-1)$. Both must vanish, yielding the
descent-swap condition $w(s_i \sigma) = t \cdot w(\sigma)$. By
Theorem~\ref{thm:rigidity}, $w(\sigma) = w(\mu) \cdot t^{\ell(\sigma)}$.
\end{proof}

\section{Discussion: implications for the bigraded lift}

The rigidity result is not a mere failure of one attempted proof. It shows
that the entire family of weight functions compatible with the
adjoint-pairing method is one-dimensional: parametrised by the global scalar
$w(\mu) \in R$. Any element of this family recovers yesterday's identity up
to a global multiplier and cannot see finer $\sigma$-dependent bigrading.

Concretely, the failure of (★) is not accidental. Even had we tried the
functional
\[
  L^{\prime}(f) \;=\; \sum_\sigma t^{\ell(\sigma)} q^{f(\sigma)}[x^\sigma] f
\]
for some sequence-valued function $f: \orb(\mu) \to \mathbb Z_{\ge 0}$, the
cancellation would demand $q^{f(s_i\sigma)} = q^{f(\sigma)}$, i.e., $f$ constant
on the orbit. Any such $q^{f(\sigma)} = q^{c}$ is a global scalar $q^c$ and
contributes nothing new.

If a genuine bigraded coset Poincar\'e identity exists, one of the following
must occur.

\begin{enumerate}
\item \textbf{Enrich the atoms.} Replace $\Aalt_\gamma$ by a bigraded
      polynomial $\Aalt^{(t,q)}_\gamma(x)$, and use yesterday's $L$ (only
      $t$-graded). If $L(\Aalt^{(t,q)}_\gamma) = \delta_{\gamma,\mu}\, F(t,q)$
      for some $F(t,q) \in \mathbb Z[t,q]$, then applying $L$ to a
      $q$-refined expansion $P^{(t,q)}_\mu = \sum c^{(t,q)}_\gamma \Aalt^{(t,q)}_\gamma$
      picks out $c^{(t,q)}_\mu = P^{(t,q)}_\mu(\dots) / F(t,q)$. But this
      requires designing both $\Aalt^{(t,q)}_\gamma$ and $P^{(t,q)}_\mu$
      coherently — no natural candidate is presently in view.
\item \textbf{Change the proof machinery.} Give up adjoint pairing and use
      a Foata-type bijection (Sec.\ \ref{sec:foata} below), or use the
      Carlsson--Chou descent basis of the Garsia--Procesi module $R_\mu$
      whose bigraded Hilbert series is Szendr\H{o}i's $A_\mu(t,q)$.
\item \textbf{Accept the negative.} The scalar identity is essentially
      tight: $c_\mu(t) = A_\mu(1,t)$, but there is no natural
      atom-decomposition-based bigraded lift.
\end{enumerate}

\subsection{Aside: what a Foata proof would look like}\label{sec:foata}

MacMahon's classical bi-Mahonian identity says $\sum_\sigma t^{\des(\sigma)} q^{\maj(\sigma)}
= \sum_\sigma t^{\des(\sigma)} q^{\inv(\sigma)}$ on rearrangements, with a
bijective proof by Foata (1968). Yesterday's identity $c_\mu(t) = \sum_\sigma
t^{\ell(\sigma)}$ can be reinterpreted as $\sum_\sigma t^{\inv(\sigma)}$
(since $\ell(\sigma) = \inv(\sigma)$ relative to $\mu$-word ordering; see
Remark~\ref{rem:inv} below). A Foata-bijective proof of the scalar identity
would first prove $c_\mu = \sum t^{\inv}$ via the atom expansion (already
done as Route B in [Clio 2026-07-28, Sec.\ 3.4] up to the atom-triangularity
step), then invoke MacMahon/Foata to conclude
$c_\mu = \sum t^{\des} q^{\maj}\big|_{t\to 1, q\to t}$. That is a
\emph{one-parameter} statement (in $t$); it does not give a genuine
$(t,q)$-bigraded refinement. In particular, it does not fill the gap
identified by Corollary~\ref{cor:no-q}.

\begin{remark}\label{rem:inv}
$\ell(\sigma) = |\{(i<j) : \sigma_i < \sigma_j\}|$ (co-inversions of the
$\mu$-word), which equals $\inv$ of the $\mu$-word with reversed alphabet
ordering. Under this identification, yesterday's $L$-functional is the
Bruhat-length weighting, and by MacMahon we have
$\sum t^{\ell(\sigma)} = A_\mu(1,t) = \binom{n}{m_1,\ldots}_t$.
\end{remark}

\section{Verification and code}

\subsection{Probe 1: falsification of (★)}
Script: \texttt{\char`\~/projects/probes/2026-07-29-bigraded-coset-poincare/test\_bigraded\_identity.py}.
Verifies that $L^{(t,q)}(\Aalt_\gamma) \ne \delta_{\gamma,\mu} t^{\des(\gamma)}
q^{\maj(\gamma)}$ for every off-diagonal $\gamma$ in each of
$\mu \in \{(2,1,0),(2,2,0),(2,2,1),(3,2,1),(2,2,0,0)\}$.
Result: 0 of 5 cases satisfy (★) (only the diagonal $\gamma=\mu$ succeeds,
trivially). Sample output for $\mu = (2,2,0)$:
\begin{verbatim}
  γ=(0,2,2), (des,maj)=(0,0): L^(t,q)(A^alt_γ) = -q*t**2 + 1     expected 0  MISS
  γ=(2,0,2), (des,maj)=(1,1): L^(t,q)(A^alt_γ) = -q*t*(q*t - 1)  expected 0  MISS
  γ=(2,2,0), (des,maj)=(1,2): L^(t,q)(A^alt_γ) = q**2*t          expected q**2*t  OK
\end{verbatim}

\subsection{Probe 2: sanity check that $\sum c_\gamma(t) L^{(t,q)}(\Aalt_\gamma) = A_\mu(t,q)$}
Script: \texttt{\char`\~/projects/probes/2026-07-29-bigraded-coset-poincare/probe2\_structure.py}.
Verified across all test cases: the weighted sum does equal $A_\mu(t,q)$,
confirming that the individual $L^{(t,q)}(\Aalt_\gamma)$ conspire correctly
--- just not diagonally.

\subsection{Probe 3: numerical confirmation of the rigidity theorem}
Script: \texttt{\char`\~/projects/probes/2026-07-29-bigraded-coset-poincare/probe3\_rigidity.py}.
Sets $w(\mu)$ symbolic, propagates $w(s_i \sigma) = t \cdot w(\sigma)$ along
descent edges of $\orb(\mu)$, and checks that the result matches
$w(\mu) \cdot t^{\ell(\sigma)}$ on every orbit element. Verified for
$\mu \in \{(2,1,0),(2,2,0),(2,2,1),(3,2,1),(2,2,0,0),(2,2,1,1)\}$: no
contradictions, exact match to $w(\mu) \cdot t^{\ell(\sigma)}$.

\section{Conclusion}

The naive bigraded lift $L^{(t,q)}$ falsifies (★), and the failure is
structural: the cancellation constraint driving yesterday's adjoint-pairing
proof forces the weight function to be $C \cdot t^{\ell(\sigma)}$, admitting
no genuine $q$-dependence per orbit element. Hence any bigraded coset
Poincar\'e identity refining $c_\mu(t) = A_\mu(1,t)$ to
$c^{(t,q)}_\mu = A_\mu(t,q)$ requires a strictly different proof route:
either bigraded atoms with a compatibly refined $P_\mu$, or a Carlsson--Chou
descent-basis argument in the Garsia--Procesi module $R_\mu$, or an entirely
new machinery. This session closes off the direct route; the next PROVE
cycle should either pivot to a descent-basis approach or return to the
chain-regime $d \ge 3$ unconditional fallback.

\end{document}
